% Resumen para CARTOGRASS 2026, Simposio Ibérico de Cartografía y Seguimiento
% de Praderas Marinas (Cartagena, 24-26 de noviembre de 2026; cartograss.eu).
% Normas: máximo 300 palabras; español, portugués o inglés; título, autores y
% afiliaciones, texto libre, hasta cinco palabras clave; sólo pósteres (A0
% vertical); plazo 30 de septiembre de 2026.
% Versión elegida entre tres redacciones (2026-09-25): arranque por el
% gradiente ecológico (decímetros de elevación = horas de emersión), método y
% validación con cifras, cierre en hábitat potencial. Números del manuscrito
% (docs/paper, ejecuciones del 23-25 de septiembre de 2026).
% Compilar:  pdflatex abstract_cartagena.tex
\documentclass[11pt,a4paper]{article}
\usepackage[utf8]{inputenc}
\usepackage[T1]{fontenc}
\usepackage[spanish,english]{babel}
\usepackage[margin=2.5cm]{geometry}
\usepackage{lmodern}
\usepackage{microtype}
\usepackage[hidelinks]{hyperref}
\pagestyle{empty}
\setlength{\parskip}{0.5em}
\setlength{\parindent}{0pt}

\begin{document}

\selectlanguage{spanish}

\begin{center}
{\large\bfseries Topografía e hidroperíodo intermareales desde Sentinel-2 con
corrección de la propagación de la marea: capa base para el hábitat potencial
de \emph{Zostera noltei}}\\[0.5em]
[Autores --- POR RELLENAR]\\
{\small [Afiliaciones --- POR RELLENAR]; correspondencia: [correo]}\\[0.3em]
{\small CARTOGRASS 2026, Cartagena, 24--26 de noviembre. Modalidad: póster.}
\end{center}

\textbf{Resumen.} En una llanura intermareal, pocos decímetros de elevación
cambian en horas el tiempo de emersión por ciclo de marea. Ese hidroperíodo
gobierna desecación y luz y ordena la zonación de \emph{Zostera noltei} en
las rías cantábricas y gallegas. Esa topografía casi nunca está medida:
levantarla es costoso y el LiDAR aéreo devuelve la superficie del agua, no el
fondo, allí donde vuela sobre la llanura inundada.

MAREA aplica el método de la línea de agua al archivo Sentinel-2 (10 m): con
la lectura mojado/seco de cada adquisición y su nivel de marea, un modelo
probit invierte la cota de cada píxel. Los productos publicados asignan un
solo nivel de agua por imagen, de un modelo oceánico o un mareógrafo; en un
estuario la onda de marea se retrasa y deforma al propagarse, y el error de
cota crece hacia el interior. MAREA estima ese retardo del propio archivo, sin
mareógrafos interiores: cada píxel actúa como mareógrafo binario y, por bandas
de 3 km desde la bocana, el retardo se obtiene por máxima verosimilitud,
monótono hacia el interior, con intervalo bootstrap.

La validación es independiente. Con pares de mareógrafos en cuatro estuarios,
los retardos estimados reproducen el orden de los medidos (MAREA/medido, min:
Ferrol 0/+4; Escalda +7/+19; Ems-Dollard +63/+51; Vlie +129/+88) y aplicarlos
reduce el error del nivel interior un 13--46\,\%. Frente a RTK-GNSS en la
ría de Villaviciosa: RMSE 0,19 m, pendiente 0,72, $r$ 0,86 (alternativa
logística publicada: 0,27 m y 0,39); frente a la batimetría Vaklodingen en
tres estuarios holandeses: 0,33--0,44 m.

Ejecutado sin supervisión en 130 celdas de 10 km de la costa cantábrica y
gallega (286 km$^2$ intermareales), entrega elevación, hidroperíodo y zonación
por píxel, con código y datos abiertos: una capa base para
cartografiar y monitorizar el hábitat potencial de \emph{Z.\ noltei}.

\textbf{Palabras clave:} hidroperíodo; línea de agua; propagación de la marea;
Sentinel-2; \emph{Zostera noltei}.

\vspace{1.2em}
\hrule
\vspace{1.2em}

\selectlanguage{english}

\begin{center}
{\large\bfseries Intertidal topography and hydroperiod from Sentinel-2 with
tide-propagation correction: a base layer for the potential habitat of
\emph{Zostera noltei}}\\[0.5em]
[Authors --- TO FILL]\\
{\small [Affiliations --- TO FILL]; corresponding author: [e-mail]}
\end{center}

\textbf{Abstract.} On an intertidal flat, a few decimetres of elevation change
the emersion time per tidal cycle by hours. That hydroperiod governs
desiccation and light and orders the zonation of \emph{Zostera noltei} in the
Cantabrian and Galician rias. This topography is almost never measured:
surveying it is costly and airborne LiDAR returns the water surface, not the
bed, wherever it flies over the flooded flat.

MAREA applies the waterline method to the Sentinel-2 archive (10 m): from the
wet/dry reading of each acquisition and its tide level, a probit model inverts
the elevation of each pixel. Published products assign a single water level
per image, from an ocean tide model or a tide gauge; inside an estuary the
tidal wave is delayed and deformed as it propagates, and the elevation error
grows landward. MAREA estimates that delay from the archive itself, without
inner tide gauges: each pixel acts as a binary tide gauge and, in 3 km bands
of distance from the mouth, the delay is obtained by maximum likelihood,
monotone landward, with a bootstrap interval.

Validation is independent. Against tide-gauge pairs in four estuaries, the
estimated delays reproduce the order of the measured ones (MAREA/measured,
min: Ferrol 0/+4; Scheldt +7/+19; Ems-Dollard +63/+51; Vlie +129/+88) and
applying them reduces the inner water-level error by 13--46\,\%. Against
RTK-GNSS in the Ría de Villaviciosa: RMSE 0.19 m, slope 0.72, $r$ 0.86
(published logistic alternative: 0.27 m and 0.39); against the Vaklodingen
bathymetry in three Dutch estuaries: 0.33--0.44 m.

Run unsupervised over 130 cells of 10 km along the Cantabrian and Galician
coast (286 km$^2$ of intertidal), it delivers per-pixel elevation, hydroperiod
and zonation, with open code and data: a base layer for mapping
and monitoring the potential habitat of \emph{Z.\ noltei}.

\textbf{Keywords:} hydroperiod; waterline method; tide propagation;
Sentinel-2; \emph{Zostera noltei}.

\end{document}
